\documentclass[12pt]{article}
\usepackage[utf8]{inputenc}
\usepackage[T1]{fontenc}
\usepackage[margin=1in]{geometry}
\usepackage{setspace}
\usepackage{xcolor}
\usepackage[colorlinks=true, urlcolor=blue!60!black, linkcolor=blue!60!black]{hyperref}
\setlength{\parindent}{2em}
\setlength{\parskip}{1.25em}
\renewcommand{\baselinestretch}{1.0}
\begin{document}
\begin{center}
{\LARGE \textit{Chae Chan Ping v. United States}, 130 U.S. 581 (1889)\par}
\vspace{0.8em}
{\large Case Brief --- Four-Page Edition\par}
\end{center}
\vspace{1em}

\section{Memory Jogger}

In \textit{The Chinese Exclusion Case}, a unanimous Court upheld the Scott Act's bar on the return of Chinese laborers holding reentry certificates, holding that a later statute overrides an earlier conflicting treaty and that Congress's power to exclude noncitizens is an inherent incident of sovereignty whose exercise is conclusive on the courts (\href{https://supreme.justia.com/cases/federal/us/130/581/}{130 U.S. 581}).

\section{Detailed Case Facts}

Chae Chan Ping, a Chinese laborer, lived lawfully in San Francisco from 1875 to 1887. The Burlingame Treaty of 1868 recognized ``the inherent and inalienable right of man to change his home and allegiance.'' After anti-Chinese agitation on the Pacific Coast, the Angell Treaty of 1880 let the United States ``regulate, limit, or suspend''---but not ``absolutely prohibit''---Chinese labor immigration, while guaranteeing that laborers already here could ``go and come of their own free will and accord.'' The Chinese Exclusion Act of 1882, as amended in 1884, suspended new labor immigration and made reentry certificates the ``only evidence permissible'' for returning residents.

Chae Chan Ping obtained a certificate before sailing for China in June 1887. While he was at sea on his return in 1888, Congress passed the Scott Act, barring all departed Chinese laborers from returning and declaring every certificate ``void and of no effect.'' On his arrival in San Francisco on October 8, 1888, the Collector refused to let him land, and he was detained aboard ship.

\section{Procedural History}

Chae Chan Ping petitioned the federal circuit court in Northern California for habeas corpus, arguing that the Scott Act violated the treaties and took a vested right in violation of the Fifth Amendment. The court denied the writ, and he appealed. The Supreme Court heard argument in March 1889 and affirmed on May 13, 1889.

\section{Judicial Votes}

Affirmed, \textbf{8--0}.
\begin{itemize}
\item \textbf{Opinion of the Court:} Justice Field, for a unanimous Court; the ninth seat was vacant after Justice Matthews's death.
\item \textbf{Concurrences and dissents:} None. (Field later dissented in \textit{Fong Yue Ting} (1893), arguing that resident noncitizens could not be deported without Fifth and Sixth Amendment protections.)
\end{itemize}

\section{Holding}

Under the Supremacy Clause, statutes and treaties stand on equal footing, so a later Act of Congress abrogates a conflicting treaty as domestic law. The power to exclude noncitizens is an inherent, inalienable incident of national sovereignty belonging to the federal government, and the legislative department's determination to exclude classes of noncitizens is conclusive upon the Judiciary. A reentry certificate confers no vested right.

\section{Analysis of Opinions}

\textbf{Later-in-time rule.} Field conceded that the Scott Act violated the treaties but held it enforceable. Relying on the \textit{Head Money Cases} and \textit{Whitney v. Robertson}, he reasoned that when a self-executing treaty and a statute conflict, ``the last expression of the sovereign will must control.'' Whether Congress had cause to abrogate a treaty, and any grievance China might press, are political questions for the other branches.

\textbf{Inherent sovereignty.} Field did not rest the exclusion power on any single enumerated clause. Drawing on \textit{The Schooner Exchange}, \textit{Cohens v. Virginia}, and \textit{Knox v. Lee}, he derived it from independent nationhood: territorial sovereignty includes the right to exclude foreigners for self-preservation, in peace as in war. In explicitly racialized terms reflecting Pacific Coast nativism, he declared that if Congress deems ``the presence of foreigners of a different race in this country, who will not assimilate with us, to be dangerous to its peace and security,'' its judgment is ``conclusive upon the judiciary.''

\textbf{Inalienability.} Sovereign powers are held in trust and cannot be bargained away. Unlike vested property rights under a treaty, a reentry certificate is a personal license ``held at the will of the government, revocable at any time, at its pleasure.''

\textbf{Later developments.} \textit{Nishimura Ekiu} (1892), \textit{Fong Yue Ting} (1893), \textit{Knauff} (1950), and \textit{Mezei} (1953) built on the decision; \textit{Wong Wing} (1896) and \textit{Kleindienst v. Mandel} (1972) set limits; and \textit{Trump v. Hawaii} (2018) applied deferential review to an entry proclamation.

\section{Examples: Future Applications}

\begin{enumerate}
\item \textbf{Same-Side (Statute Controls): Visa-Waiver Treaty.} Congress bars all nationals of a treaty partner and voids their visas. A later statute abrogates the treaty, and a visa is a revocable license.
\item \textbf{Same-Side (Statute Controls): Tax Treaty Override.} A later statute taxes interest a treaty exempted; ``the last expression of the sovereign will must control,'' and breach is for diplomacy.
\item \textbf{Opposite-Side (Not Shielded): Hard Labor Without a Jury.} Imposing infamous punishment on noncitizens inside the country requires full criminal protections under \textit{Wong Wing}.
\item \textbf{Opposite-Side (No Deference): Proclamation Contrary to Statute.} A presidential ban on asylum applications contrary to 8 U.S.C.\ \S\ 1158 puts executive power at its ``lowest ebb'' under \textit{Youngstown}; the plenary power belongs to Congress.
\item \textbf{Fence-Sitter: Religious Exclusion Statute.} Congress bars immigrants of a named religion. \textit{Chae Chan Ping} deferred even to racial exclusions, but \textit{Trump v. Hawaii} stressed facial neutrality.
\end{enumerate}

\section{Critique}

\textbf{Extratextual sovereignty versus enumerated powers.} The decision departs from \textit{McCulloch}'s premise of a government of enumerated powers. Deriving unenumerated ``inherent'' sovereignty from international law risks a constitutional black hole, a domain untethered from Article I limits and Bill of Rights constraints.

\textbf{Racial animus and unchecked deference.} Field's reliance on racial unassimilability shows the peril of treating exclusion decisions as ``conclusive upon the judiciary.'' Defenders say the doctrine has been cleansed of racial justifications while keeping needed foreign-affairs deference; critics answer that \textit{Trump v. Hawaii} shows how easily plenary power still shields discriminatory action.

\section{Key Quotations}

\begin{quote}
``For local interests the several states of the Union exist, but for national purposes, embracing our relations with foreign nations, we are but one people, one nation, one power. ... If, therefore, the government of the United States, through its legislative department, considers the presence of foreigners of a different race in this country, who will not assimilate with us, to be dangerous to its peace and security, their exclusion is not to be stayed because at the time there are no actual hostilities with the nation of which the foreigners are subjects. ... In both cases its determination is conclusive upon the judiciary.''
\end{quote}
\noindent --- Justice Field, Opinion of the Court, 130 U.S.\ at 606.

\begin{quote}
``By the constitution, laws made in pursuance thereof, and treaties made under the authority of the United States, are both declared to be the supreme law of the land, and no paramount authority is given to one over the other. ... In either case the last expression of the sovereign will must control.''
\end{quote}
\noindent --- Justice Field, Opinion of the Court, 130 U.S.\ at 600.

\begin{quote}
``That the government of the United States, through the action of the legislative department, can exclude aliens from its territory is a proposition which we do not think open to controversy. Jurisdiction over its own territory to that extent is an incident of every independent nation. It is a part of its independence.''
\end{quote}
\noindent --- Justice Field, Opinion of the Court, 130 U.S.\ at 603--604.

\end{document}
